\documentclass[10pt,a4paper]{amsart}
\usepackage[margin=19mm]{geometry}
\usepackage{amsmath,amssymb,mathtools}
\usepackage[scr=rsfs,cal=euler]{mathalfa}
\usepackage{xfrac}
\usepackage[unicode,hidelinks,bookmarks=false]{hyperref}
\newcommand{\ie}{i.e.\ }
\newcommand{\cf}{cf.\ }
\newcommand{\resp}{respectively\ }
\newcommand{\Subsection}{\subsection}
\providecommand{\nobreakdash}{\nobreak\hbox{-}}
\setlength{\emergencystretch}{2em}
\multlinegap0pt
\usepackage{amssymb}
\usepackage{mathtools}
\usepackage{tikz}
\usepackage{xcolor}
\usepackage{tcolorbox}
\usepackage[shortlabels]{enumitem}
\newtheorem{proposition}{Proposition}
\title{Geodesic Quantum $f$-Divergences}
\author{\'Angela Capel, Pablo Costa Rico}
\date{Source: arXiv:2608.15833v1 (CC BY 4.0)}
\begin{document}
\maketitle

\noindent\emph{Source and licence.} Geodesic Quantum $f$-Divergences. Version arXiv:2608.15833v1 (CC BY 4.0), distributed by arXiv under the Creative Commons Attribution 4.0 International licence, http://creativecommons.org/licenses/by/4.0/, as recorded in the arXiv licence metadata for this exact version and shown on the abstract page https://arxiv.org/abs/2608.15833v1. Full licence notice: \texttt{licenses/arxiv-2608.15833-CC-BY-4.0.txt}.\par\smallskip

\noindent\textit{Editorial scope.} Counted unit: the article's proposition rw:prop:alpha-z-intersections (source0.tex lines 1819--1860). The article's complete original proof of it is delivered with it (source0.tex lines 1862--2040, one \texttt{proof} environment of 179 lines). No mathematics is written, repaired or re-derived: the statement and the proof are the article's own lines, in the article's own notation, and no retained line is substituted or re-flowed.  Retained uncounted as the source-local context the proof needs: lines 1634--1644; lines 1806--1814.  Nothing was omitted from any retained span, so the package carries no omission record.  No retained line contains a citation, so the wrapper carries no bibliography and no bibliography entry.  No retained line contains a figure environment, an image-inclusion command or drawing code, so no image is delivered and the package declares no asset dependency.  Editorial formatting only, in addition: an AMS \texttt{amsart} preamble that restates the article's own declarations and macros used by the retained lines, namely the environments \texttt{proposition} on separate counters, together with the standard packages those lines need.  The retained environments print in this document as Proposition 1 in order of appearance, whereas the article prints the same items with the numbering of its own full document.  The complete native source of the article as distributed in the arXiv e-print for this exact version is bundled unchanged as \texttt{sources/207787/source0.tex}.
\begin{align}
 D_\alpha^0(\rho\Vert\sigma)
 &=\frac{1}{\alpha-1}
   \log\operatorname{Tr}(\rho^\alpha\sigma^{1-\alpha}),
 \label{rw:eq:petz-renyi-endpoint}\\
 D_\alpha^1(\rho\Vert\sigma)
 &=\frac{1}{\alpha-1}
   \log\operatorname{Tr}\!\left[
    \sigma(\sigma^{-1/2}\rho\sigma^{-1/2})^\alpha\right],
 \label{rw:eq:geometric-renyi-endpoint}
\end{align}

\begin{equation}
 Q_{\alpha,t}
 =\operatorname{Tr}\!\left[
  \rho^{\alpha-(\alpha-1)t}
  \bigl(\rho^{t/2}\sigma^{-1}\rho^{t/2}\bigr)^{\alpha-1}
 \right],
 \qquad t\in\mathbb R.
 \label{rw:eq:power-functional-push-through}
\end{equation}

\begin{proposition}[Exact universal intersections for real $t$]
\label{rw:prop:alpha-z-intersections}
We assume $\dim\mathcal H\geq2$ and fix $\alpha>0$, $\alpha\neq1$,
$z>0$, and $t\in\mathbb R$.  Then
\begin{equation}
 D_\alpha^t(\rho\Vert\sigma)
 =D_{\alpha,z}(\rho\Vert\sigma)
 \quad\text{for every invertible pair }(\rho,\sigma)
 \label{rw:eq:alpha-z-universal-equality}
\end{equation}
if and only if at least one of the following conditions holds:
\begin{equation}
 \begin{aligned}
 &t=0,\quad z=1;\qquad
 &&\alpha=2,\quad z=1;\\
 &t=\frac{\alpha}{\alpha-1},\quad
   z=\lvert\alpha-1\rvert;\qquad
 &&(\alpha,t,z)=(3,3,1).
 \end{aligned}
 \label{rw:eq:alpha-z-intersection-locus}
\end{equation}
The alternatives form a union and may overlap.  Equivalently, the
universal identities are the Petz line $D_\alpha^0=D_{\alpha,1}$,
the full real order-two fibre
\begin{equation}
 D_2^t=D_{2,1}
 =\log\operatorname{Tr}(\rho^2\sigma^{-1}),
 \qquad t\in\mathbb R,
 \label{rw:eq:alpha-z-order-two-fibre}
\end{equation}
the extrapolated branch
\begin{equation}
 D_\alpha^{\,\alpha/(\alpha-1)}
 =D_{\alpha,\lvert\alpha-1\rvert},
 \label{rw:eq:alpha-z-extrapolated-intersection}
\end{equation}
and the isolated cubic coincidence
\begin{equation}
 D_3^3=D_{3,1}.
 \label{rw:eq:alpha-z-cubic-intersection}
\end{equation}
\end{proposition}

\begin{proof}
The Petz identity follows from
\eqref{rw:eq:petz-renyi-endpoint}.  We put $\beta:=\alpha-1$.  From
\eqref{rw:eq:power-functional-push-through}, cyclicity gives
\begin{equation}
 Q_{2,t}=\operatorname{Tr}(\rho^2\sigma^{-1}),
 \qquad t\in\mathbb R,
 \label{rw:eq:real-t-quadratic-identity}
\end{equation}
which proves the quadratic fibre.  At
$t_\star:=\alpha/\beta$, the exterior power of $\rho$ in
\eqref{rw:eq:power-functional-push-through} disappears.  If
$\beta>0$, the resulting trace is the $(\alpha,z)$ functional at
$z=\beta$; if $\beta<0$, we use
\begin{equation}
 \bigl(\rho^{t_\star/2}\sigma^{-1}
       \rho^{t_\star/2}\bigr)^{-1}
 =\rho^{-t_\star/2}\sigma\rho^{-t_\star/2}
 \label{rw:eq:negative-branch-inversion}
\end{equation}
to obtain the $(\alpha,z)$ functional at $z=-\beta$.  This proves
\eqref{rw:eq:alpha-z-extrapolated-intersection}.  Finally, for the
real-$t$ continuation,
\begin{align}
 Q_{3,t}
 &=\operatorname{Tr}\!\left(
   \sigma^{-1}\rho^{3-t}\sigma^{-1}\rho^t\right)\notag\\
 &=\left\|\rho^{(3-t)/2}\sigma^{-1}\rho^{t/2}\right\|_2^2
 =\left\|\rho^{t/2}\sigma^{-1}\rho^{(3-t)/2}\right\|_2^2
 =Q_{3,3-t},
 \label{rw:eq:cubic-reflection}
\end{align}
Thus $Q_{3,3}=Q_{3,0}=\operatorname{Tr}(\rho^3\sigma^{-2})$, which
is precisely the $(\alpha,z)$ functional at $(\alpha,z)=(3,1)$.

We prove exhaustiveness.  A two-dimensional calculation is enough.
Indeed, if an invertible qubit pair $(\rho_2,\sigma_2)$ gives a mismatch,
then for any invertible state $\tau$ on the complementary subspace and
$0<p<1$, we set
\begin{equation}
 \rho=p\rho_2\oplus(1-p)\tau,
 \qquad
 \sigma=p\sigma_2\oplus(1-p)\tau.
 \label{rw:eq:real-t-qubit-embedding}
\end{equation}
Each power functional is $p$ times its qubit value plus $1-p$, so
the mismatch persists in every dimension.  We take
\begin{equation}
 \rho_x:=\frac{1}{1+x}
 \begin{pmatrix}1&0\\0&x\end{pmatrix},
 \qquad
 S:=\sigma^{-1}
 =\begin{pmatrix}a&c\\ \overline c&b\end{pmatrix}>0,
 \label{rw:eq:real-t-rank-one-witness}
\end{equation}
and $P=\lvert0\rangle\!\langle0\rvert$.  We set
$m_r:=\langle0\rvert S^r\lvert0\rangle$.  As $x\downarrow0$,
\begin{equation}
 Q_{\alpha,z}(\rho_x\Vert\sigma)
 \longrightarrow m_{\beta/z}^{\,z}.
 \label{rw:eq:real-t-alpha-z-pure-limit}
\end{equation}
To record the relevant spectral calculus, for $t>0$ and up to a
scalar tending to one,
\begin{equation}
 K_x:=\rho_x^{t/2}S\rho_x^{t/2}
 =\begin{pmatrix}
   a&cx^{t/2}\\ \overline c x^{t/2}&bx^t
  \end{pmatrix}.
 \label{rw:eq:real-t-spectral-block}
\end{equation}
Its eigenvalues and the lower diagonal entry of the spectral
projection $E_+$ of the larger eigenvalue satisfy
\begin{align}
 \lambda_+(K_x)=a+O(x^t), \qquad
 \lambda_-(K_x)=\frac{ab-\lvert c\rvert^2}{a}x^t+O(x^{2t}), \qquad
 \langle1\rvert E_+\lvert1\rangle
 =\frac{\lvert c\rvert^2}{a^2}x^t+O(x^{2t}).
 \label{rw:eq:real-t-spectral-asymptotics}
\end{align}
In \eqref{rw:eq:power-functional-push-through}, the latter term is
weighted by $x^{\alpha-\beta t}$, producing the competing exponent
$\kappa:=\alpha+t(2-\alpha)$.  At $\kappa=0$ its extra contribution
is $\lvert c\rvert^2a^{\beta-2}$, so together with $a^\beta$ it
equals $m_1^{\beta-2}m_2$.  Applying the same calculation to
$K_x^{-1}=\rho_x^{-t/2}\sigma\rho_x^{-t/2}$ yields the negative-$t$
cases.  Consequently, for $\alpha\neq2$,
\begin{equation}
 Q_{\alpha,t}(\rho_x\Vert\sigma)\longrightarrow
 \begin{cases}
  m_\beta,&t=0,\\[1mm]
  m_{-1}^{-\beta},&t<0,\ 0<\alpha<2,\\[1mm]
  m_1^\beta,&t>0,\ \kappa>0,\\[1mm]
  m_1^{\beta-2}m_2,&t>0,\ \kappa=0.
 \end{cases}
 \label{rw:eq:real-t-pure-limits}
\end{equation}
For $c\neq0$ the same quantity diverges when $t<0$, $\alpha>2$,
or when $t>0$, $\kappa<0$.

We let $M_r:=m_r^{1/r}$ be the scalar power mean of the spectral values
of $S$ in the vector state $P$.  Since
$m_{\beta/z}^z=M_{\beta/z}^{\beta}$, strict monotonicity of power
means in \eqref{rw:eq:real-t-alpha-z-pure-limit} and \eqref{rw:eq:real-t-pure-limits}
forces
\begin{equation}
 \begin{array}{c|c}
  t=0&z=1\\
  t<0&\beta<0,\ z=-\beta\\
  t>0,\ \kappa>0&\beta>0,\ z=\beta.
 \end{array}
 \label{rw:eq:real-t-power-mean-candidates}
\end{equation}
At $\alpha=2$, \eqref{rw:eq:real-t-quadratic-identity} and the first
line of the same power-mean argument force $z=1$.

It remains to analyze the positive-$t$ threshold $\kappa=0$, which
can occur only for $\alpha>2$ and
$t=\alpha/(\alpha-2)$.  We put $p:=\beta/z$ and take
$S_u=\left(\begin{smallmatrix}1&u\\u&1\end{smallmatrix}\right)$,
$\lvert u\rvert<1$.  Normalizing $S_u^{-1}$ to a state only
multiplies both sides below by the same positive scalar.  Universal
equality would imply
\begin{equation}
 \left(
  \frac{(1+u)^p+(1-u)^p}{2}
 \right)^{\beta/p}=1+u^2.
 \label{rw:eq:real-t-threshold-identity}
\end{equation}
The coefficient of $u^2$ gives $\beta(p-1)=2$; after substitution,
the coefficient of $u^4$ gives
$(\beta-2)(2\beta+1)=0$.  Since $\beta>1$, this forces
$\beta=p=2$.  Hence the threshold contributes exactly
$(\alpha,t,z)=(3,3,1)$.

We finally test the two regular candidates with
$z=\lvert\beta\rvert$.  We put
\begin{equation}
 \delta:=t_\star-t,\qquad
 A:=\rho^\delta,\qquad
 B:=\rho^{t/2}\sigma^{-1}\rho^{t/2}.
 \label{rw:eq:real-t-ALT-variables}
\end{equation}
Then
\begin{equation}
 Q_{\alpha,t}=\operatorname{Tr}(A^\beta B^\beta),
 \qquad
 \left.Q_{\alpha,z}\right|_{z=\lvert\beta\rvert}
 =\operatorname{Tr}(A^{1/2}BA^{1/2})^\beta.
 \label{rw:eq:alpha-z-ALT-comparison}
\end{equation}
If $\delta=0$, this is precisely the extrapolated branch.  If
$\delta\neq0$, arbitrary positive definite $A$ and $B$ can be
realized up to positive scalar factors by taking
$\rho\propto A^{1/\delta}$ and
$\sigma\propto\rho^{t/2}B^{-1}\rho^{t/2}$.
For $\beta>0$, we choose $A=\operatorname{diag}(a,1)$, $a>0$,
$a\neq1$, and
$B_\varepsilon=(1-\varepsilon)P_++\varepsilon P_-$, where
$P_\pm$ project onto $2^{-1/2}(1,\pm1)^{\mathsf T}$.  The two traces
in \eqref{rw:eq:alpha-z-ALT-comparison} tend to
\begin{equation}
 \frac{a^\beta+1}{2}
 \quad\text{and}\quad
 \left(\frac{a+1}{2}\right)^\beta,
 \label{rw:eq:alpha-z-strict-witness}
\end{equation}
which are unequal unless $\beta=1$.  If $\beta=-q<0$, multiplying
the two traces by $\varepsilon^q$ instead gives the limits
\begin{equation}
 \frac{a^{-q}+1}{2}
 \quad\text{and}\quad
 \left(\frac{a+1}{2a}\right)^q,
 \label{rw:eq:alpha-z-negative-strict-witness}
\end{equation}
which are strictly unequal for $0<q<1$ by concavity of $x^q$.
Thus $\delta=0$, except when $\beta=1$; that exception is exactly the
quadratic fibre already found.  This proves that the list is complete.
\end{proof}

\end{document}
